% SPDX-License-Identifier: MIT
% Copyright (c) 2026 Evgeny Baulin
% Optimization Methods, Seminar 08: one-page cheat sheet (English).
% Generated by checks/08. Constrained Optimality and Discrete Modelling.ipynb from web/08/js/content-cheatsheet.js and content/students.json. Do not edit by hand.
% Compile with: latexmk -xelatex -interaction=nonstopmode -halt-on-error Cheatsheet_en.tex
\RequirePackage{fix-cm}
\documentclass[10pt,a4paper]{article}

\usepackage{amsmath,amssymb}
\usepackage{fontspec}
\setmainfont{cmunrm.otf}[BoldFont=cmunbx.otf, ItalicFont=cmunti.otf, BoldItalicFont=cmunbi.otf]
\setsansfont{cmunss.otf}[BoldFont=cmunsx.otf, ItalicFont=cmunsi.otf, BoldItalicFont=cmunso.otf]
\setmonofont{cmuntt.otf}[BoldFont=cmuntb.otf, ItalicFont=cmunit.otf, BoldItalicFont=cmuntx.otf]
\usepackage{polyglossia}
\setdefaultlanguage{english}

\usepackage[a4paper,left=8mm,right=8mm,top=7mm,bottom=11mm,footskip=5.5mm]{geometry}
\usepackage{xcolor}
\usepackage{multicol}
\usepackage{enumitem}
\usepackage{ragged2e}
\usepackage{fancyhdr}
\usepackage[colorlinks=true,linkcolor=hseblue,urlcolor=hseblue]{hyperref}

\definecolor{hseblue}{RGB}{16,45,105}
\definecolor{hseaccent}{RGB}{198,40,40}
\definecolor{hsegray}{RGB}{108,114,126}
\definecolor{hselight}{RGB}{240,243,248}

\hypersetup{pdftitle={Seminar 08 cheat sheet: Constrained Optimality and Discrete Modelling}, pdfauthor={Evgeny Baulin}}

\pagestyle{fancy}
\fancyhf{}
\renewcommand{\headrulewidth}{0pt}
\fancyfoot[L]{\scriptsize\color{hsegray}\textcopyright\ 2026 Evgeny Baulin, \href{https://t.me/tarakan_tuc}{t.me/tarakan\_tuc}}
\fancyfoot[R]{\scriptsize\color{hsegray}MIT License}

\setlength{\parindent}{0pt}
\setlength{\parskip}{0pt}
\setlength{\columnsep}{4.5mm}
\setlength{\columnseprule}{0.3pt}
\renewcommand{\columnseprulecolor}{\color{hselight}}
\setlength{\emergencystretch}{2em}
\setlength{\RaggedRightRightskip}{0pt plus 4em}
\binoppenalty=2000
\relpenalty=500
\setlength{\abovedisplayskip}{2pt}
\setlength{\belowdisplayskip}{2pt}
\setlist[itemize]{leftmargin=7pt, labelsep=3pt, itemsep=1.2pt, topsep=1pt, parsep=0pt, partopsep=0pt,
  label={\color{hseaccent}\scriptsize\textbullet}}

\makeatletter
\newcommand{\cssec}[1]{\par\addvspace{3.5pt}\noindent
  \colorbox{hselight}{\parbox{\dimexpr\linewidth-2\fboxsep}{\RaggedRight\hyphenpenalty=10000\color{hseblue}\bfseries\strut #1}}%
  \par\nobreak\vspace{1pt}\@afterheading}
\makeatother
\clubpenalty=10000
\widowpenalty=10000

\begin{document}
\thispagestyle{fancy}

\noindent
\begin{minipage}[b]{0.6\textwidth}
  \RaggedRight\hyphenpenalty=10000
  {\color{hseaccent}\bfseries\small Seminar 08 cheat sheet}\\[1pt]
  {\color{hseblue}\bfseries\Large Constrained Optimality\\ and Discrete Modelling}
\end{minipage}\hfill
\begin{minipage}[b]{0.39\textwidth}
  \raggedleft\footnotesize
  {\bfseries Optimization Methods}\\
  HSE University, Faculty of Computer Science, 2026--2027\\
  Course lecturer: Andrey Ignatov\\
  Author: Evgeny Baulin, Telegram \href{https://t.me/tarakan_tuc}{@tarakan\_tuc}
\end{minipage}

\vspace{2pt}
{\color{hseblue}\rule{\textwidth}{0.8pt}}
\vspace{-4pt}

\fontsize{9}{11.4}\selectfont
\begin{multicols}{3}
\RaggedRight


% ---- kkt
\cssec{KKT conditions}
\begin{itemize}
  \item $L=f+\sum\lambda_ig_i+\sum\nu_jh_j$, $g_i\le0$, $\lambda_i\ge0$, $\nu_j$ free.
  \item Stationarity $\nabla_xL=0$, feasibility, $\lambda\ge0$, $\lambda_ig_i=0$.
  \item Necessary under LICQ (or, for a convex problem, a strictly feasible point); sufficient when the problem is convex.
  \item Slack $\Rightarrow\lambda_i=0$; tight does not imply $\lambda_i>0$.
\end{itemize}

% ---- cases
\cssec{Active-set recipe}
\begin{itemize}
  \item Pick $A$: $g_i=0$ on $A$, $\lambda_i=0$ off $A$; solve.
  \item Reject if some $g_i>0$ or $\lambda_i<0$.
  \item Convex: the survivor is global. Nonconvex: compare all survivors.
\end{itemize}

% ---- dual
\cssec{Duality and prices}
\begin{itemize}
  \item $q=\inf_xL$ is concave; $q(\lambda,\nu)\le f(x)$ for feasible $x$, $\lambda\ge0$.
  \item Convex + strictly feasible point: $d^*=p^*$, dual optimum attained.
  \item $g_i\le b_i$: $\partial p^*/\partial b_i=-\lambda_i^*$. For $\ge$ the sign flips.
  \item $f(x)-q(\lambda)\ge0$ certifies how far a feasible $x$ can be from optimal.
\end{itemize}

% ---- milp
\cssec{Integer models}
\begin{itemize}
  \item Decisions $\to$ variables, rules $\to$ constraints, goal $\to$ objective.
  \item Relaxation: max gives an upper bound, min a lower bound; rounding is not a method.
  \item $x\le Mz$ switches a fixed cost; $z_A\le z_B$ is if A then B; $\sum z_i\le k$ caps a choice.
  \item Take $M$ as small as valid: same integer solutions, tighter relaxation.
\end{itemize}

\end{multicols}
\end{document}
